\documentclass[12pt]{article}
\usepackage[utf8]{inputenc}
\usepackage[T1]{fontenc}
\usepackage[margin=1in]{geometry}
\usepackage{setspace}
\usepackage{xcolor}
\usepackage[colorlinks=true, urlcolor=blue!60!black, linkcolor=blue!60!black]{hyperref}
\setlength{\parindent}{2em}
\setlength{\parskip}{1.25em}
\renewcommand{\baselinestretch}{1.0}
\begin{document}
\begin{center}
{\LARGE \textit{Buck v. Bell}, 274 U.S. 200 (1927)\par}
\vspace{0.8em}
{\large Case Brief\par}
\end{center}
\vspace{1em}

\section{Memory Jogger}

In \href{https://www.courtlistener.com/opinion/101076/buck-v-bell/}{\textit{Buck v. Bell}}, the Supreme Court upheld Virginia's compulsory eugenic sterilization statute against Fourteenth Amendment due process and equal protection challenges, analogizing forced salpingectomy to compulsory vaccination and declaring that ``[t]hree generations of imbeciles are enough.''

\section{Detailed Case Facts}

On March 20, 1924, the Virginia General Assembly enacted a compulsory sterilization statute drafted by state senator Aubrey E. Strode (\href{https://encyclopediavirginia.org/entries/buck-v-bell-1927/}{Encyclopedia Virginia}; \href{https://supreme.justia.com/cases/federal/us/274/200/}{274 U.S. at 205}). The statute recited that the health of the patient and the welfare of society could be promoted by sterilizing ``mental defectives'' via vasectomy (for males) or salpingectomy (for females) without serious pain or substantial danger to life, asserting that heredity played an important role in transmitting insanity and imbecility and that sterilized inmates could be safely discharged to become self-supporting (\href{https://supreme.justia.com/cases/federal/us/274/200/}{274 U.S. at 205--206}). To insulate the law from constitutional attack, the act required a multi-step procedure: a verified petition by the institution's superintendent to its special board of directors, notice to the inmate and a court-appointed guardian (and parents if a minor), a hearing with written evidence that the inmate could attend, and rights of appeal to the county circuit court and the Virginia Supreme Court of Appeals (\href{https://supreme.justia.com/cases/federal/us/274/200/}{274 U.S. at 206--207}).

The test case centered on Carrie Elizabeth Buck, an eighteen-year-old white woman at the time of her late-1924 trial, who had been committed to the Virginia State Colony for Epileptics and Feeble-Minded near Lynchburg (\href{https://supreme.justia.com/cases/federal/us/274/200/}{274 U.S. at 205}; \href{https://encyclopediavirginia.org/entries/buck-v-bell-1927/}{Encyclopedia Virginia}). In the formal judicial record accepted by the Supreme Court, Carrie Buck was described as a ``feeble minded white woman,'' the daughter of a ``feeble minded mother'' (Emma Buck) confined in the same institution, and the mother of an ``illegitimate feeble minded child'' (Vivian Buck) (\href{https://supreme.justia.com/cases/federal/us/274/200/}{274 U.S. at 205}). Subsequent historical investigation by legal scholar Paul A. Lombardo and others revealed a starkly different reality: Carrie had been raised by foster parents, became pregnant after being raped by their nephew, and was committed to conceal the family scandal; neither Carrie nor her daughter Vivian---who later made the elementary school honor roll before dying at age eight---was intellectually disabled (\href{https://encyclopediavirginia.org/entries/buck-v-bell-1927/}{Encyclopedia Virginia}; \href{https://www.law.cornell.edu/supremecourt/text/18-483}{Cornell LII, \textit{Box v. Planned Parenthood}}).

Historically and politically, the litigation unfolded at the height of the American eugenics movement, which blended Mendelian heredity theories with Progressive-era faith in scientific social engineering (\href{https://encyclopediavirginia.org/entries/buck-v-bell-1927/}{Encyclopedia Virginia}). Colony Superintendent Albert S. Priddy deliberately selected Carrie Buck as a test case to validate Virginia's 1924 statute before launching broader sterilizations. Buck's counsel, Irving P. Whitehead, was a former colony board member and sterilization supporter who failed to cross-examine state witnesses aggressively, called no witnesses of his own, and conceded Carrie's ``feeblemindedness'' in his brief (\href{https://encyclopediavirginia.org/entries/buck-v-bell-1927/}{Encyclopedia Virginia}). Buck challenged the sterilization order under the Due Process and Equal Protection Clauses of the Fourteenth Amendment.

\section{Procedural History}

On September 10, 1924, the special board of directors of the Virginia State Colony approved Superintendent Albert S. Priddy's petition to sterilize Carrie Buck (\href{https://encyclopediavirginia.org/entries/buck-v-bell-1927/}{Encyclopedia Virginia}). At Priddy's instigation, Buck's court-appointed guardian, Robert G. Shelton, appealed the order to the Circuit Court of Amherst County (\textit{Buck v. Priddy}). Following a hearing in November 1924, Judge Bennett T. Gordon upheld the sterilization order in February 1925 (issuing a written judgment on April 13, 1925), finding that Buck was feebleminded and ``the probable potential parent of socially inadequate offspring, likewise afflicted'' (\href{https://supreme.justia.com/cases/federal/us/274/200/}{274 U.S. at 207}; \href{https://encyclopediavirginia.org/entries/buck-v-bell-1927/}{Encyclopedia Virginia}).

After Priddy's death in January 1925, the colony's new superintendent, Dr. John H. Bell, was substituted as respondent. On November 12, 1925, the Supreme Court of Appeals of Virginia affirmed the circuit court's judgment (143 Va. 310) (\href{https://supreme.justia.com/cases/federal/us/274/200/}{274 U.S. at 205}; \href{https://encyclopediavirginia.org/entries/buck-v-bell-1927/}{Encyclopedia Virginia}). Buck sought review in the U.S. Supreme Court by writ of error. Following oral argument on April 22, 1927, the Supreme Court affirmed the judgment of the Virginia Supreme Court of Appeals on May 2, 1927 (\href{https://supreme.justia.com/cases/federal/us/274/200/}{274 U.S. at 200, 208}). Dr. Bell performed the salpingectomy on Carrie Buck on October 19, 1927 (\href{https://encyclopediavirginia.org/entries/buck-v-bell-1927/}{Encyclopedia Virginia}).

\section{Judicial Votes}

The decision was 8--1 (\href{https://supreme.justia.com/cases/federal/us/274/200/}{Justia}; \href{https://en.wikipedia.org/wiki/Buck_v._Bell}{Wikipedia}):
\begin{itemize}
  \item \textbf{Majority (8):} Justice Oliver Wendell Holmes Jr. (author), joined by Chief Justice William Howard Taft and Justices Willis Van Devanter, James Clark McReynolds, Louis D. Brandeis, George Sutherland, Edward Terry Sanford, and Harlan Fiske Stone.
  \item \textbf{Concurrences (0):} None.
  \item \textbf{Dissent (1):} Justice Pierce Butler (dissented without written opinion).
\end{itemize}

\section{Holding}

\textbf{Judgment.} Affirmed. The order directing the superintendent of the Virginia State Colony to perform a salpingectomy on Carrie Buck was upheld (\href{https://supreme.justia.com/cases/federal/us/274/200/}{274 U.S. at 208}).

\textbf{Rule Announced.} A state statute authorizing the compulsory sexual sterilization of inmates in state-supported institutions who are adjudicated, following thorough procedural safeguards, to be afflicted with hereditary forms of insanity or imbecility does not violate the Due Process Clause of the Fourteenth Amendment, as it falls within the state's police power to protect public welfare. Furthermore, confining the statute's application to inmates of state institutions rather than extending it to similarly situated individuals outside those institutions does not violate the Equal Protection Clause, because a state may address a perceived social problem incrementally as its means allow (\href{https://supreme.justia.com/cases/federal/us/274/200/}{274 U.S. at 207--208}).

\textbf{Holding vs. Dicta.} The holding is that Virginia's compulsory sterilization statute, as applied to an institutionalized person adjudicated under statutory hearing and appeal procedures to carry hereditary mental defects, satisfies both Fourteenth Amendment Due Process and Equal Protection. Justice Holmes's sweeping rhetoric comparing sterilization to military conscription, warning against ``being swamped with incompetence,'' and proclaiming that ``[t]hree generations of imbeciles are enough'' reflects his polemical rationale rather than a precise doctrinal threshold---the statute did not legally require three generations of impairment, and the Court did not rule on sterilizing non-institutionalized persons or using sterilization as criminal punishment (\href{https://supreme.justia.com/cases/federal/us/274/200/}{274 U.S. at 207--208}).

\section{Analysis of Opinions}

\subsection{Majority Opinion (Justice Holmes)}

Justice Holmes's remarkably brief opinion (spanning barely three pages in the United States Reports) disposed of Buck's procedural due process, substantive due process, and equal protection claims in rapid succession.

First, Holmes summarily rejected any procedural due process defect. After recounting the Virginia statute's notice, guardian-appointment, board-hearing, written-record, and dual-level judicial appeal provisions, he concluded that ``the rights of the patient are most carefully considered'' and that, because every statutory step was followed after months of observation, Buck unquestionably received procedural due process (\href{https://supreme.justia.com/cases/federal/us/274/200/}{274 U.S. at 206--207}).

Second, turning to substantive due process, Holmes deferred completely to the Virginia legislature's eugenic premises and the state courts' factual findings that Buck was the probable parent of ``socially inadequate offspring'' whose sterilization would promote both her welfare and society's (\href{https://supreme.justia.com/cases/federal/us/274/200/}{274 U.S. at 207}). Framing the issue as a utilitarian trade-off between individual bodily autonomy and collective state strength, Holmes analogized forced sterilization to wartime conscription---reasoning that if the state can demand the lives of its ``best citizens,'' it may demand ``these lesser sacrifices'' from those ``who already sap the strength of the State.'' He then invoked \textit{Jacobson v. Massachusetts}, 197 U.S. 11 (1905), asserting that ``[t]he principle that sustains compulsory vaccination is broad enough to cover cutting the Fallopian tubes,'' and capped the paragraph with his notorious epigram: ``Three generations of imbeciles are enough'' (\href{https://supreme.justia.com/cases/federal/us/274/200/}{274 U.S. at 207}).

Third, Holmes dismissed the equal protection challenge---that Virginia unconstitutionally singled out institutionalized inmates while leaving the ``multitudes outside'' untouched---as ``the usual last resort of constitutional arguments'' (\href{https://supreme.justia.com/cases/federal/us/274/200/}{274 U.S. at 208}). He held that the Equal Protection Clause permits incremental legislative action: the law ``does all that is needed when it does all that it can, indicates a policy, applies it to all within the lines, and seeks to bring within the lines all similarly situated so far and so fast as its means allow.'' Moreover, Holmes added an ironic twist: because sterilizing inmates enabled their safe release and freed asylum beds for others, the statute actually advanced equality by bringing more people within its reach (\href{https://supreme.justia.com/cases/federal/us/274/200/}{274 U.S. at 208}).

\subsection{Dissenting Opinion (Justice Butler)}

Justice Pierce Butler dissented without writing an opinion, leaving only the bare notation ``MR. JUSTICE BUTLER dissents'' (\href{https://supreme.justia.com/cases/federal/us/274/200/}{274 U.S. at 208}). Historians have often attributed Butler's lone dissent to his devout Roman Catholicism (which opposed eugenic sterilization on natural-law grounds) and his broader commitment to substantive individual liberty against aggressive state police-power intrusions. Because Butler left no written reasoning, however, his dissent provided no doctrinal counter-framework for lower courts to invoke.

\section{Examples: Future Applications}

The following five hypotheticals illustrate the reach and limits of \textit{Buck v. Bell}'s doctrinal reasoning (taken strictly on its own 1927 terms, prior to later constitutional developments):

\begin{enumerate}
  \item \textbf{Same-Side Application (Compulsory Sterilization of Institutionalized Male Inmates):} A state enacts a statute modeled on Virginia's 1924 law authorizing compulsory vasectomies for male residents of state psychiatric hospitals found, after a formal board hearing with counsel and judicial appeal, to suffer from hereditary severe cognitive impairment. \textit{Why:} This falls squarely within \textit{Buck v. Bell}: the state provides robust procedural safeguards, asserts a public-welfare rationale grounded in hereditary transmission, and applies the measure uniformly within state institutions.

  \item \textbf{Same-Side Application (Incremental Public Health Mandate):} A state legislature requires mandatory vaccination against a contagious disease only for residents of state-run residential facilities and public university dormitories, omitting private households due to limited enforcement resources. \textit{Why:} Under \textit{Buck}'s equal protection holding, a state does not violate the Fourteenth Amendment by addressing a public-health problem incrementally (``within the lines'' of state-managed populations) as its administrative means allow.

  \item \textbf{Opposite-Side Application (Summary Sterilization Without Hearing or Appeal):} A state hospital superintendent orders the immediate surgical sterilization of an incoming psychiatric patient based solely on an intake nurse's unverified notes, with no prior notice to the patient or a guardian, no evidentiary hearing, and no opportunity for judicial review. \textit{Why:} Even under \textit{Buck v. Bell}, Holmes emphasized that Virginia's statute satisfied procedural due process precisely because of its ``very careful provisions''---notice, guardian representation, a written-record hearing, and two levels of judicial appeal. Summary sterilization fails \textit{Buck}'s procedural floor.

  \item \textbf{Opposite-Side Application (Arbitrary Punitive Sterilization of Selected Felons):} A state statute mandates compulsory sterilization for persons convicted three times of larceny (such as chicken theft) while exempting persons convicted three times of embezzlement, without any individualized finding of hereditary medical defect. \textit{Why:} As the Supreme Court later held in \textit{Skinner v. Oklahoma}, 316 U.S. 535 (1942) (\href{https://encyclopediavirginia.org/entries/buck-v-bell-1927/}{Encyclopedia Virginia}), distinguishing arbitrarily between functionally identical crimes for punitive sterilization violates Equal Protection and falls outside \textit{Buck}'s non-punitive, institutional-health rationale.

  \item \textbf{Fence-Sitter Application (Compulsory Reversible Contraception for Non-Institutionalized Welfare Recipients):} A state statute provides full administrative hearings and judicial appeals before requiring long-acting, reversible contraceptive implants for unmarried parents diagnosed with hereditary intellectual disabilities who live in the community and receive state financial assistance. \textit{Why:} This is a close question under \textit{Buck}'s own logic. On one side, the intervention is less physically permanent than salpingectomy, includes Virginia-style procedural safeguards, and invokes the state's fiscal interest in preventing ``socially inadequate'' dependency. On the other side, it extends beyond institutionalized wards whose release depends on sterilization, directly implicating the bodily autonomy of free citizens outside state custody.
\end{enumerate}

\section{Critique}

\textbf{Scholarly Criticism and Logical Weaknesses.} \textit{Buck v. Bell} sits firmly within the American constitutional ``anti-canon,'' alongside \textit{Dred Scott}, \textit{Plessy}, and \textit{Korematsu} (\href{https://digitalcommons.lib.uconn.edu/cgi/viewcontent.cgi?article=1495&context=law_review}{Connecticut Law Review}). Its analytical flaws are glaring. First, Holmes's reliance on \textit{Jacobson v. Massachusetts} committed a massive false equivalence: a smallpox vaccination is a minor, temporary preventive measure against an acute, communicable epidemic, whereas surgical sterilization is a permanent, invasive destruction of reproductive capacity aimed at non-communicable, poorly understood hereditary traits. Second, Holmes uncritically accepted a sham factual record produced by collusive litigation---treating pseudoscientific eugenic assertions about Carrie, Emma, and Vivian Buck as settled truths (\href{https://encyclopediavirginia.org/entries/buck-v-bell-1927/}{Encyclopedia Virginia}; \href{https://www.law.cornell.edu/supremecourt/text/18-483}{Cornell LII}). Third, Holmes's equal protection analysis inverted the clause's protective purpose: treating the state's eventual goal of sterilizing even more people as a virtue that ``more nearly reached'' equality reduced equal protection to a rubber stamp. The decision catalyzed over 60,000 involuntary sterilizations in the United States and was later cited by Nazi defendants at the Nuremberg Trials (\href{https://www.law.cornell.edu/supremecourt/text/18-483}{Cornell LII}; \href{https://encyclopediavirginia.org/entries/buck-v-bell-1927/}{Encyclopedia Virginia}). Although \textit{Skinner v. Oklahoma} (1942) gutted its rationale, the Supreme Court has never formally overruled \textit{Buck v. Bell} (\href{https://en.wikipedia.org/wiki/Buck_v._Bell}{Wikipedia}).

\textbf{Competing Values and Premises.} At its core, \textit{Buck v. Bell} pitted extreme majoritarian utilitarianism and technocratic deference against fundamental human dignity and bodily integrity. Holmes treated citizens as instruments of the collective state---valuable when they strengthen it, expendable when they ``sap'' it---and elevated judicial restraint into an abdication of constitutional review even when the state permanently mutilates an individual's body.

\textbf{Progressive Perspective.} Modern progressive, feminist, and disability-rights scholars condemn \textit{Buck v. Bell} as a chilling example of how state power, class bias, gender policing, and ableism masquerade as neutral science (\href{https://www.cambridge.org/core/books/feminist-judgments-reproductive-justice-rewritten/buck-v-bell-274-us-200-1927/AFDA01C302E0A630A86607743D21ED3F}{Cambridge Univ. Press, \textit{Feminist Judgments}}). Carrie Buck's commitment punished a poor young woman for an out-of-wedlock pregnancy resulting from sexual assault, illustrating how eugenic categories targeted marginalized women and people with disabilities. Progressive jurisprudence also underscores the historical irony that \textit{Buck} was authored by Holmes and joined by Louis Brandeis---heroes of early-twentieth-century Progressive judicial restraint---demonstrating that uncritical deference to legislative ``expertise'' without heightened scrutiny for fundamental reproductive and equality rights invites grave abuses.

\textbf{Originalist and Textualist Perspective.} Originalists and textualists likewise repudiate \textit{Buck v. Bell}, though on distinct grounds. First, under an original-meaning reading of the Fourteenth Amendment's Privileges or Immunities Clause and Due Process Clause, bodily integrity and freedom from arbitrary physical mutilation by the state are classic common-law rights deeply rooted in Anglo-American liberty that no valid ``police power'' extends to destroying absent criminal conviction or imminent contagion. Second, conservative jurists such as Justice Clarence Thomas---who recounted the sordid history of \textit{Buck v. Bell} at length in his concurrence in \href{https://www.supremecourt.gov/opinions/18pdf/18-483_3d9g.pdf}{\textit{Box v. Planned Parenthood of Indiana and Kentucky, Inc.}}, 587 U.S. 490 (2019)---cite \textit{Buck} as a cautionary monument to elite Progressive social engineering and the dangers of treating human beings as ``undesirable populations'' subject to state-sanctioned reproductive selection (\href{https://www.law.cornell.edu/supremecourt/text/18-483}{Cornell LII}). For originalists, \textit{Buck} proves that neither uncabined judicial policymaking nor blind deference to fashionable scientific consensus is a substitute for enforcing the Constitution's textually anchored protections of individual liberty and equality.

\section{Key Quotations}

All quotations below are verified verbatim from Justice Holmes's majority opinion in \href{https://www.courtlistener.com/opinion/101076/buck-v-bell/}{\textit{Buck v. Bell}}, 274 U.S. 200 (1927) (\href{https://supreme.justia.com/cases/federal/us/274/200/}{Justia}):

\begin{quote}
``There can be no doubt that, so far as procedure is concerned, the rights of the patient are most carefully considered, and, as every step in this case was taken in scrupulous compliance with the statute and after months of observation, there is no doubt that, in that respect, the plaintiff in error has had due process of law.'' \hfill (274 U.S. at 207)
\end{quote}

\begin{quote}
``We have seen more than once that the public welfare may call upon the best citizens for their lives. It would be strange if it could not call upon those who already sap the strength of the State for these lesser sacrifices, often not felt to be such by those concerned, in order to prevent our being swamped with incompetence.'' \hfill (274 U.S. at 207)
\end{quote}

\begin{quote}
``It is better for all the world if, instead of waiting to execute degenerate offspring for crime or to let them starve for their imbecility, society can prevent those who are manifestly unfit from continuing their kind. The principle that sustains compulsory vaccination is broad enough to cover cutting the Fallopian tubes. Jacobson v. Massachusetts, 197 U. S. 11. Three generations of imbeciles are enough.'' \hfill (274 U.S. at 207)
\end{quote}

\begin{quote}
``It is the usual last resort of constitutional arguments to point out shortcomings of this sort. But the answer is that the law does all that is needed when it does all that it can, indicates a policy, applies it to all within the lines, and seeks to bring within the lines all similarly situated so far and so fast as its means allow.'' \hfill (274 U.S. at 208)
\end{quote}

\end{document}
